\documentclass[11pt]{article}
\usepackage[a4paper,margin=25mm]{geometry}
\usepackage{amsmath,amssymb,amsthm,booktabs,hyperref}
\setlength{\emergencystretch}{2em}
\newtheorem{theorem}{Theorem}
\newcommand{\F}{\mathbb F}\newcommand{\PSL}{\operatorname{PSL}}\newcommand{\PSU}{\operatorname{PSU}}
\title{Chiral constructions and regular subgroups in dimension three}
\author{Principal synthesis for Leemans--Vandenschrick}\date{}
\begin{document}\maketitle
We recover the three principal theorems of \emph{On chiral polytopes
having a group $\PSL(3,q)$ as automorphism group}, in its fixed
25-page author edition. The constructions and their original matrices
are due to Leemans and Vandenschrick. The accompanying papers give
complete new proofs, distinguish their classical inputs, and record
several corrections to intermediate formulas and tables. The final
unitary examples in Section~7 are supplied with explicit generators.
Earlier computational statements cited there remain attributed prior
results, with their evidence boundary stated below.

The common ideas are the actual translation kernels of small projective
subgroups, elementary matrix generation, and the action of an
involution's centre and axis. A subgroup order does not by itself prove
an intersection property: the constructions use the actual common
subgroups, and the finite unitary examples check every indexed
intersection. A separate centralizer calculation excludes reversing
automorphisms of the whole ambient group.

\section*{The three principal conclusions}
Let $q=p^f$, and let $d=\gcd(3,q-1)$. We use the standard string
rotation correspondence for abstract regular and chiral polytopes.
\begin{theorem}\label{principal}
The following assertions hold.
\begin{enumerate}
\item For every $q\ge5$, the three matrices in Theorem~1.1 of the source,
with a primitive element $x\in\F_q^\times$, generate a self-dual chiral
polytope with automorphism group $\PSL_3(q)$ and type
\[
                   [q-1,\,2(q-1)/d,\,q-1].
\]
The adjacent subgroups are the full set stabilizers of two distinct
lines and two distinct points, respectively, of order
$2q^2(q-1)^2/d$.
\item For every prime $p\equiv1\pmod6$, each of the four markings
specified in Theorem~1.2 of the source gives a chiral polytope with
automorphism group $\PSL_3(p)$. Their types are $[3,6,6,3]$ for the
positive sign and $[3,6,3,3]$ for the negative sign. Including duals
there are six distinct oriented marking classes in this construction,
which give three unbased polytope classes after identifying the two
opposite base-flag orbits. The same matrices over $\F_{p^2}$, with
$p\equiv5\pmod6$, give directly regular examples: $\PSL_3(p^2)$ is
the rotation group and its extension by Frobenius is the full group.
\item For every $q\ge2$, neither $\PSL_3(q)$ nor $\PSU_3(q)$ is the
automorphism group of a chiral polytope of rank greater than five.
\end{enumerate}
\end{theorem}
\begin{proof}
The first accompanying paper, \emph{Two stabilizers and elementary
matrices in dimension three}, retains the literal source matrices and
duality. Its finite-field span argument generates both full translation
radicals, hence all six elementary root subgroups. This proves ambient
generation without a maximal-subgroup classification. The stabilizer
intersection gives the rotation intersection property. Its final
semilinear calculation includes the inverse-transpose coset and proves
chirality, while direct multiplication gives the type and self-duality.
This proves the first assertion for every prime power in the stated
range, rather than only the fields used in supplementary controls.

\emph{A parabolic subgroup and a conic in the rank-five construction}
proves the exact translation module
of the first adjacent subgroup and the restricted parabolic containing
it. The two signs are treated separately. For the negative sign a
literal symmetric-square representation gives the conic end; the
positive sign uses the stated point-pair duality. Both actual middle
intersections and every recursive rotation intersection are proved.
The common centralizer of three prefix involutions is scalar, excluding
all reversing automorphisms over the prime field. Frobenius reverses
the quadratic-field markings. The comparison of both cube roots and
both signs, including duals, gives exactly the oriented-class statement
above. This proves the second assertion with its orientation and full
regular-group qualifications.

For the third assertion, the paper \emph{The two regular ends of a chiral marking}
first derives the two directly regular ends from the actual indexed
rotation intersection property. The acquired regular subgroup
classification bounds their ranks. A rotation of order two gives a
proper normal direct product, so it cannot occur in the relevant simple
ambient groups. The elementary unitary argument for $\PSU_3(2)$ pays
the nonsimple boundary separately. It remains to exclude rank six.
In even characteristic the two modules and their actual centres and
axes force a pair of distant rotations not to commute. In odd
characteristic the Heisenberg normal form gives the same contradiction
by an explicit commutator. The remaining directly regular simplex case
has rotation group a quotient of $A_7$; group orders exclude it in all
ambient fields of characteristic five. The complete argument is uniform
and proves the claimed bound, including all its small-field premises.
\end{proof}

\section*{The auxiliary classifications and source tables}
The source's Section~5 is recovered by four accompanying components.
\emph{Small regular subgroups and their translation kernels} proves the
rank-two order classification and the complete odd-characteristic
alternatives in rank three. \emph{Elations and two translation modules}
proves the even-characteristic normal forms, the exact field and
intersection conditions, and the regular rank bound.
\emph{Four homologies in a projective plane} gives the odd rank-four
classification, including its degenerate alternative. The conic branch
reuses the already proved projective-line classification: the exceptional
markings are those for $\PSL_2(11)$, $\PSL_2(19)$ and $\operatorname{PGL}_2(5)$.
Their exact presentations and checks remain acquired inputs.

These proofs use Dickson's subgroup types and the classical
three-dimensional conic representation as attributed inputs. They do
not count those classifications as new results. In the odd affine
branch, writing $K=\F_p(a)$ identifies the actual kernel as $K^2$;
the complements are literal translates, so their intersections can be
calculated before appealing to the string end criterion. In even
characteristic the corresponding two modules determine precisely when
the marking has the intersection property. This accounts for all group
alternatives and types used in Tables~2--3. The earlier rank-three
projective-line classifications cited by those tables remain attributed
inputs, with the actual string and intersection hypotheses retained.
No automatic passage from an arbitrary hypermap marking is used.

The separate calculation \emph{The parameter restrictions in Table 1}
pays all nine signed parameter rows of that auxiliary table. It computes
the three characteristic polynomials and eliminates $t$ from $t^2=5$
and the exact middle-order equations. The listed primes are necessary
parameter restrictions; the actual square root must belong to the
ambient field and satisfy the displayed equation. Neither the table
nor its heading is used as an independent intersection certificate.

The source corrections have precise limits. The first elementary-matrix
word on physical page~6 needs an $x^{-1}$ factor in one coordinate;
the generation paragraph on page~7 interchanges its end-subgroup names.
In Section~4 the conic sign requires its own intersection argument; the
adjacent translation kernel over a quadratic field has order $p^2$,
and the regular analogue has an additional Frobenius automorphism.
The six-class sentence is qualified by the oriented-marking convention.

On page~11 a regular unipotent has order $p$ in odd characteristic,
not the printed bound $2(2,p)$. The final dihedral divisibility bound
survives. On page~17 the direct product is $D_{2p}^2$, not $D_{2p}$;
the final characteristic-polynomial coefficient is $a'$, not $b$.
The point/line statements are qualified by duality. The simplex
alternative needed in Section~6 occurs in characteristic five, allowing
embeddings into extension fields, rather than only the single ambient
field of order five.

Finally, the generic affine row of Table~3 needs a type correction.
If $\lambda+\lambda^{-1}=2a$, put $h=|\lambda|$ and
$h'=|{-\lambda}|$. Its actual types are
\[
                       [h,h'],\quad[h,2p],\quad[h',2p],
\]
not just the printed $[k,4]$ and $[k,k]$. For example $p=5,a=2$
gives three C-markings in the group of order $300$, with types
$[3,6]$, $[3,10]$ and $[6,10]$. The uniform proof and full finite
checks agree. This repairs a row without changing the principal
constructions or the rank obstruction.

\section*{Section 7: positive examples and retained prior results}
\emph{Three explicit unitary examples in rank four} recovers the three
positive $\PSU_3(17)$ examples reported in the source, with types
$[3,7,3]$, $[3,8,3]$ and $[3,9,3]$. Its matrices are defined by three
Hermitian centres over $\F_{17}[\beta]$, $\beta^2=3$. Fresh exact
permutation calculations prove full ambient generation and all $768$
indexed intersections. A written calculation of the full semilinear
centralizer excludes every reversing automorphism. This is a positive
existence statement; no absence of other types is inferred from the
source's sentence or from the exploratory search.

The negative lists for $\PSL_3(q)$ with $q\le4$ and for
$\PSU_3(q)$ with $q\in\{2,3,4,5,7,9,13,16\}$ are explicitly
attributed in the source to earlier computational work. The three
$\PSU_3(11)$ pairs, of types $[5,4,5]$ twice and $[4,5,4]$ once,
are likewise attributed there to data collected for Buekenhout,
Leemans and Tranchida. We retain these as prior source-attributed
results, not as our independently repeated censuses. The available
edition of that earlier paper provides timings and group data but not
those output tuples; the project's accepted algorithmic recovery does
not manufacture the missing execution receipts. This evidence boundary
is explicit in the retained-input register.

The same paragraph's negative list for $\PSL_2(q)$ contains one
incorrect entry: $q=81$. The acquired projective-line theorem gives a
chiral example of type $[40,40,40]$. Indeed its pointed construction
applies with $p=3$, $d=4$ and $k=40$: $k$ is even and divides
$(81-1)/2$, while none of $3^e-1,3^e+1$ for $1\le e<4$ is divisible
by $40$. This corrects the quoted list by an already proved theorem.
The remaining entries of that list satisfy its negative criteria.

The previously known absence of chiral polyhedra for $\PSL_3(q)$ and
$\PSU_3(q)$ is the theorem attributed to Breda d'Azevedo and Catalano.
Together with Theorem~\ref{principal}(3), it gives the source's
conclusion that a unitary chiral example can have only rank four or
five. The other introductory classification statements are retained
with their earlier authorship; the already acquired project results
supply exact reuse links where available.

The historical question about unitary groups for $q>17$ is recorded
without a new claim about its present status. The quoted composite-degree
$\PSL_2$ question is settled by the previously acquired coprime-subfield
construction in the projective-line paper. Its full proof and earlier
attribution are retained.
\end{document}
